\documentclass[11pt]{article}
\usepackage[margin=1in]{geometry}
\usepackage{amsmath,amssymb,amsthm}
\usepackage{booktabs}
\usepackage{hyperref}
\usepackage{lmodern}
\usepackage{enumitem}

\newtheorem{theorem}{Theorem}
\newtheorem{definition}[theorem]{Definition}
\newtheorem{proposition}[theorem]{Proposition}
\newtheorem{corollary}[theorem]{Corollary}
\newtheorem{lemma}[theorem]{Lemma}
\theoremstyle{remark}
\newtheorem{remark}[theorem]{Remark}

\newcommand{\phig}{\varphi}
\newcommand{\mustar}{\mu_\star}
\newcommand{\Epass}{E_{\mathrm{pass}}}
\newcommand{\Ecoh}{E_{\mathrm{coh}}}
\newcommand{\Bpow}{B_{\mathrm{pow}}}
\newcommand{\Etot}{E_{\mathrm{total}}}
\newcommand{\Rdel}{R_\Delta}

\title{A $\phig$-Ladder Prediction for Neutrino Mass Ordering\\
and the Dimensionless Splitting Ratio\\[6pt]
\large A Falsifiable Prediction from the Recognition Science
Mass Backbone}

\author{Jonathan Washburn\\
Recognition Physics Institute\\
Austin, Texas, USA\\
\texttt{jon@recognitionphysics.org}}

\date{March 2026}

\begin{document}
\maketitle

\begin{abstract}
The Recognition Science mass backbone places neutrinos on deep
negative rungs of the $\phig$-ladder ($\phig=(1+\sqrt{5})/2$),
with fractional rung positions forced by octave compatibility
and the edge-partition structure of the three-cube~$Q_3$.
Given the neutral-sector quarter-step structure, the framework
makes three predictions that are independent of the SI
calibration seam:
\begin{enumerate}[nosep]
  \item \textbf{Mass ordering}: normal ($m_1<m_2<m_3$), forced
    by the monotonicity of the $\phig$-ladder.
  \item \textbf{Squared-mass ratio}: $m_3^2/m_2^2 = \phig^7
    \approx 29.03$.
  \item \textbf{Splitting ratio}:
    $\Rdel := \Delta m^2_{31}/\Delta m^2_{21}
    = (\phig^{11}-1)/(\phig^4-1) \approx 33.82$.
\end{enumerate}
The splitting ratio is the sharpest near-term test.  It is
dimensionless, depends only on the golden ratio and derived
cube integers ($\Epass=11$, $2^{D-1}=4$), and carries no
dependence on the absolute mass scale or the SI anchor.
The current NuFIT~5.3 best fit for normal ordering gives
$\Rdel^{\mathrm{exp}} = 31.6 \pm 1.5$, placing the prediction
at approximately $1.5\sigma$ tension.  If future oscillation
data push the measured $\Rdel$ beyond $3\sigma$ from~$33.82$,
the framework's neutrino rung assignment is refuted.

One structural input remains at the frontier: the $-1/4$
quarter-phase offset that selects the neutrino rung triple
from the deep-ladder window.  This paper frames the prediction
conditionally on that offset rather than claiming it as a
closed first-principles derivation.
\end{abstract}

%=============================================================================
\section{Introduction}
%=============================================================================

The neutrino mass ordering---whether the lightest neutrino
mass eigenstate is $\nu_1$ (normal ordering) or $\nu_3$
(inverted ordering)---is one of the major open questions in
particle physics.  Current oscillation experiments prefer
normal ordering but have not established it decisively.
The absolute neutrino mass scale remains unknown, bounded
from above by cosmological constraints ($\Sigma m_\nu < 0.12$~eV)
and tritium endpoint measurements.

A framework that predicts the mass ordering, the splitting
ratio, and the absolute mass scale from a discrete combinatorial
structure would be either a powerful tool or a conspicuous
target for falsification.  This paper provides such predictions
from the Recognition Science mass backbone.

The backbone paper (Paper~I~\cite{backbone}) established the
fermion mass law at the common anchor scale
$\mustar=182.201$~GeV, with sector yardsticks, integer rungs,
and the charge-band correction all derived from the three-cube
$Q_3$.  The lepton paper (Paper~II~\cite{leptons}) demonstrated
charged-lepton mass ratios from $Q_3$ integers with no free
parameters.  The CKM paper (Paper~III~\cite{CKM}) showed how
the backbone connects to Standard Model mixing observables.

The present paper applies the backbone to the neutrino sector.
The key simplification is that neutrinos are electrically
neutral ($Q=0$), so the charge-band correction vanishes
identically: $\mathrm{gap}(0)=0$.  The neutrino mass spectrum
reduces to a pure $\phig$-ladder, and every mass ratio is a
$\phig$-power of a rung difference.

This simplification makes the neutrino sector the cleanest
testing ground for the framework.  The dimensionless splitting
ratio $\Rdel$ depends on nothing but $\phig$ and the rung
differences, which are derived from cube integers.  It is
independent of the absolute mass scale, the sector yardstick,
the SI calibration seam, and even the electroweak VEV.

\subsection{Scope and framing}

This paper is a \emph{prediction paper}, not a full-derivation
paper.  It states the neutrino predictions that follow from
the neutral-sector quarter-step structure and identifies the
one structural input that remains at the frontier.  The
predictions are conditional on that input; they are sharp,
falsifiable, and honest about what is assumed versus derived.

What belongs in this paper: the neutral-sector simplification,
the quarter-step lattice, the edge-partition rung spacings, the
absolute baseline derivation, and the three seam-free predictions
(ordering, $\phig^7$ ratio, splitting ratio).

What does not belong: the full RS ontology, the charged-lepton
derivation, quark masses, or the anchor/units argument beyond a
brief citation of Paper~I.

%=============================================================================
\section{Notation and prerequisites}
\label{sec:notation}
%=============================================================================

We use the backbone conventions of Paper~I~\cite{backbone}.
The cube integers are $V=8$, $E=\Etot=12$, $F=6$, $A=1$,
$\Epass=E-A=11$, $W=\Epass+F=17$.  The golden ratio is
$\phig=(1+\sqrt{5})/2$, forced by the Recognition Composition
Law.  The octave period is $T_{\min}=2^D=8$ for spatial
dimension $D=3$.

The mass law at the anchor scale $\mustar$ (Paper~I,
Eq.~(1)) is
\begin{equation}\label{eq:masslaw}
  m_f^{\mathrm{RS}}(\mustar)
  = A_{s(f)}\,\phig^{\,r_f - 8 + \mathrm{gap}(Z_f)},
\end{equation}
where $A_s$ is the sector yardstick and
$\mathrm{gap}(Z)=\log_\phig(1+Z/\phig)$ is the charge-band
correction.

%=============================================================================
\section{Neutral-sector simplification}
\label{sec:neutral}
%=============================================================================

\begin{proposition}[Neutral-sector mass law]
\label{prop:neutral}
For neutrinos ($Q_\nu=0$), the mass law reduces to
\begin{equation}\label{eq:nu_mass}
  m_{\nu_i} = A_{\mathrm{EW}}\,\phig^{\,r_i - 8},
\end{equation}
with no charge-band correction.
\end{proposition}

\begin{proof}
$Q_\nu=0$ implies $\tilde{Q}=6Q=0$, hence
$Z_\nu = \tilde{Q}^2 + \tilde{Q}^4 = 0$, and therefore
$\mathrm{gap}(0) = \log_\phig(1+0/\phig) = \log_\phig 1 = 0$.
The neutrino sector yardstick is $A_{\mathrm{EW}}$
(electroweak sector: $\Bpow=+1$, $r_0=55$), giving
$A_{\mathrm{EW}} = 2\,\Ecoh\,\phig^{55} = 2\phig^{50}$
in framework-native units.
\end{proof}

\begin{corollary}[Pure $\phig$-power ratios]
\label{cor:ratios}
All neutrino mass ratios are pure $\phig$-powers of rung
differences:
\begin{equation}\label{eq:nu_ratio}
  \frac{m_{\nu_i}}{m_{\nu_j}} = \phig^{\,r_i - r_j}.
\end{equation}
The sector yardstick, the octave offset, and the SI calibration
all cancel in the ratio.
\end{corollary}

This cancellation is the reason the neutrino sector provides
the cleanest falsification targets.  The predictions depend on
nothing but $\phig$ and the rung triple.

%=============================================================================
\section{The quarter-step lattice}
\label{sec:quarter}
%=============================================================================

Charged fermions sit on integer rungs of the $\phig$-ladder.
The neutrino mass splittings are far smaller than the
charged-fermion hierarchy, requiring finer resolution.  The
question is: what is the minimal sub-lattice of $\mathbb{Z}$
that is compatible with the octave structure?

\begin{theorem}[Quarter-step resolution]
\label{thm:quarter}
The minimal refinement of the integer-rung lattice that is
compatible with the 8-tick octave period and accommodates a
$-1/4$ phase offset is the quarter-step lattice
$r\in\tfrac{1}{4}\mathbb{Z}$.
\end{theorem}

\begin{proof}
A sub-lattice $r\in\frac{1}{n}\mathbb{Z}$ is
\emph{octave-compatible} if $n$ divides $T_{\min}=8$.  The
divisors of~$8$ are $\{1,2,4,8\}$.  To accommodate a phase
offset of $-1/4$, the lattice must resolve quarter-integer
positions: $n\geq 4$.  The minimal such divisor is $n=4$.
\end{proof}

\begin{remark}[Frontier status of the phase offset]
\label{rem:frontier}
The $-1/4$ phase offset is the single structural input in the
neutrino sector that remains at the frontier.  A plausible
geometric origin is the $C_4$ rotational symmetry of each face
of $Q_3$ (each face is a square, providing a natural fourfold
subdivision of the integer rung).  However, this paper does not
claim a closed first-principles derivation of the phase value.
The predictions below are conditional on the $-1/4$ offset.
If the offset were instead $-1/8$ or $-3/8$, the rung triple
and all derived predictions would change.  This conditionality
is kept visible rather than hidden.
\end{remark}

%=============================================================================
\section{Rung spacings from edge partition}
\label{sec:spacings}
%=============================================================================

The three neutrino masses have two structurally independent
rung spacings: $\Delta_{21}:=r_2-r_1$ and
$\Delta_{32}:=r_3-r_2$.

\begin{theorem}[Rung spacings]
\label{thm:spacings}
The two spacings are derived from the directional partition of
the passive edges of~$Q_3$:
\begin{equation}\label{eq:spacings}
  \Delta_{21} = \frac{2^{D-1}}{2} = \frac{4}{2} = 2,
  \qquad
  \Delta_{32} = \frac{\Epass - 2^{D-1}}{2} = \frac{7}{2}.
\end{equation}
The total span is
$\Delta_{31} = \Delta_{21}+\Delta_{32} = \Epass/2 = 11/2$.
\end{theorem}

\noindent\textit{Geometric origin.}\quad
Neutrinos, being electrically neutral, couple only to the 1D
edge network of $Q_3$, not to its 2-faces.  The total rung
span of the neutrino triplet is therefore half the passive edge
count: $\Delta_{31}=\Epass/2=11/2$.  The passive edges partition
into two groups:
\begin{itemize}[nosep]
  \item \textbf{Axial group}: $2^{D-1}=4$ edges (those along
    one coordinate axis direction).
  \item \textbf{Remaining group}: $\Epass-2^{D-1}=11-4=7$ edges.
\end{itemize}
Each group maps to a rung spacing at half-resolution (edge count
divided by~2), giving the two spacings.  The halving factor
reflects the quarter-step lattice: the edge-only sub-lattice
operates at $\frac{1}{2}\mathbb{Z}$ resolution, using half the
passive edge count.

In quarter-rung numerator coordinates:
$\Delta_{21}=8/4$ and $\Delta_{32}=14/4$, both elements of
$\frac{1}{4}\mathbb{Z}$.  The total
$\Delta_{31}=22/4=11/2=\Epass/2$ exhausts the passive edge
budget.

%=============================================================================
\section{Absolute baseline rung}
\label{sec:baseline}
%=============================================================================

The spacings fix the relative positions.  The absolute position
requires one additional input: the baseline rung $r_3$ of the
heaviest neutrino.

\begin{theorem}[Neutrino baseline from hypercube integers]
\label{thm:baseline}
The integer part of the atmospheric neutrino rung is
\begin{equation}\label{eq:baseline}
  r_{\nu_3}^{\mathrm{int}}
  = -(V+E+F+\Epass+W)
  = -(8+12+6+11+17) = -54.
\end{equation}
This is the negative sum of all five non-trivial hypercube
integers.  The active edge count $A=1$ is excluded because it
enters the charged-sector rung formula separately.
\end{theorem}

The integer baseline $-54$ is formally verified in Lean~4
(\texttt{Masses.BaselineDerivation.neutrino\_baseline\_eq}).

\begin{theorem}[Unique rung triple]
\label{thm:triple}
Under the deep-ladder atmospheric confinement (the quarter-rung
numerator $4r_3$ lies in the integer interval $(-220,-216)$) and
the $-1/4$ phase offset ($4r_3+1\equiv 0\pmod{4}$), the rung
triple is uniquely determined:
\begin{equation}\label{eq:triple}
  (r_1,\,r_2,\,r_3)
  = \left(-\frac{239}{4},\;-\frac{231}{4},\;-\frac{217}{4}\right).
\end{equation}
\end{theorem}

\begin{proof}
The integers in $(-220,-216)$ are $\{-219,-218,-217\}$.  The
phase condition $4r_3+1\equiv 0\pmod{4}$ selects $4r_3=-217$
uniquely (since $-217+1=-216$ and $-216\bmod 4=0$, while
$-219+1=-218\bmod 4=2$ and $-218+1=-217\bmod 4=3$, both
failing).

Back-propagation through the fixed spacings:
$r_2 = r_3 - 7/2 = -217/4 - 14/4 = -231/4$,\;
$r_1 = r_2 - 2 = -231/4 - 8/4 = -239/4$.
\end{proof}

\begin{remark}[Phase selection and normal ordering]
Among the three deep-window candidates, only the $-1/4$ phase
($4r_3=-217$) produces normal ordering $m_1<m_2<m_3$.  The
other two candidates produce inverted ordering or degenerate
masses inconsistent with oscillation data.  The phase $-1/4$ is
therefore the unique deep-window quarter-step choice consistent
with the observed mass ordering.  It remains a frontier item
(no deeper cube-geometric derivation is available), but it is
selected by data, not freely assumed.
\end{remark}

%=============================================================================
\section{The three seam-free predictions}
\label{sec:predictions}
%=============================================================================

All three predictions below depend only on $\phig$ and the rung
differences.  They are independent of the absolute mass scale,
the sector yardstick, the SI calibration anchor $\tau_0$, and
any scheme or scale choice in the Standard Model transport layer.

\subsection{Prediction~1: Normal mass ordering}

\begin{theorem}[Normal ordering is forced]
\label{thm:ordering}
The rung triple $r_1 < r_2 < r_3$ together with $\phig > 1$
forces
\begin{equation}\label{eq:ordering}
  m_1 < m_2 < m_3
  \qquad\textup{(normal ordering)}.
\end{equation}
\end{theorem}

\begin{proof}
$m_i = C\,\phig^{r_i}$ with $C>0$ and $\phig>1$ is strictly
increasing in~$r$.  Since
$-239/4 < -231/4 < -217/4$, the ordering follows.
\end{proof}

Normal ordering is not a choice in this framework.  It is a
structural consequence.  If future experiments decisively
establish inverted ordering ($m_3 < m_1$), the rung
triple~\eqref{eq:triple} is refuted.

\subsection{Prediction~2: Squared-mass ratio $\phig^7$}

\begin{theorem}[$\phig^7$ atmospheric-to-solar ratio]
\label{thm:phi7}
\begin{equation}\label{eq:phi7}
  \frac{m_3^2}{m_2^2}
  = \phig^{\,2(r_3-r_2)}
  = \phig^{\,2\times 7/2}
  = \phig^7
  \approx 29.03.
\end{equation}
\end{theorem}

The exponent~7 is the number of remaining passive edges
($\Epass-2^{D-1}=11-4=7$): the same cube integer that
controls the atmospheric--solar rung gap.

This ratio is testable once absolute neutrino mass information
becomes available from direct kinematic measurements or
cosmological surveys.

\subsection{Prediction~3: The dimensionless splitting ratio}

\begin{theorem}[Splitting ratio]
\label{thm:splitting}
The ratio of mass-squared splittings is
\begin{equation}\label{eq:splitting}
  \boxed{
  \Rdel
  := \frac{\Delta m^2_{31}}{\Delta m^2_{21}}
  = \frac{\phig^{2\Delta_{31}}-1}{\phig^{2\Delta_{21}}-1}
  = \frac{\phig^{11}-1}{\phig^{4}-1}
  \approx 33.82.}
\end{equation}
\end{theorem}

\begin{proof}
Writing $m_i = C\,\phig^{r_i}$ for a common pre-factor $C$:
\begin{align*}
  \Delta m^2_{31}
  &= m_3^2 - m_1^2
  = C^2\phig^{2r_1}(\phig^{2\Delta_{31}}-1), \\
  \Delta m^2_{21}
  &= m_2^2 - m_1^2
  = C^2\phig^{2r_1}(\phig^{2\Delta_{21}}-1).
\end{align*}
The common factor $C^2\phig^{2r_1}$ cancels:
\[
  \Rdel
  = \frac{\phig^{2\times 11/2}-1}{\phig^{2\times 2}-1}
  = \frac{\phig^{11}-1}{\phig^{4}-1}.
\]
Numerically: $\phig^{11}\approx 199.0$,
$\phig^{4}\approx 6.854$, hence
$\Rdel\approx 198.0/5.854\approx 33.82$.
\end{proof}

The exponents $11=\Epass$ and $4=2^{D-1}$ are the same cube
integers that enter the edge sub-partition
(Theorem~\ref{thm:spacings}).  The splitting ratio is therefore
a function of the passive edge count and the cube dimension
alone.

%=============================================================================
\section{Comparison with oscillation data}
\label{sec:data}
%=============================================================================

Table~\ref{tab:seam_free} compares the three seam-free
predictions with current experimental constraints.

\begin{table}[h]
\centering
\caption{Seam-free predictions vs.\ current data.  These
predictions depend only on $\phig$ and derived cube integers;
they carry no SI-seam dependence.}
\label{tab:seam_free}
\begin{tabular}{lccc}
\toprule
Prediction & RS value & Experimental & Status \\
\midrule
Mass ordering & Normal & Preferred (NuFIT~5.3) & Consistent \\
$m_3^2/m_2^2$ & $\phig^7\approx 29.03$ & Not yet measured
  & Awaiting data \\
$\Rdel$ & $33.82$ & $31.6\pm 1.5$ (NuFIT~5.3, NO)
  & $\approx 1.5\sigma$ tension \\
\bottomrule
\end{tabular}
\end{table}

The splitting ratio tension is the strongest current pressure
on the framework.  It is reported transparently here because
honest claim hygiene requires it.

\begin{remark}[Structural interpretation of the tension]
The predicted $\Rdel=33.82$ corresponds to $\Delta_{32}=7/2$.
Inverting the splitting ratio formula with the experimental
central value $\Rdel^{\mathrm{exp}}=31.6$ and holding
$\Delta_{21}=2$ fixed gives $\Delta_{32}\approx 3.36$.  The
nearest quarter-rung values are $13/4=3.25$ and $7/2=3.50$.
Neither matches $3.36$ exactly.  The $1.5\sigma$ tension is
not resolvable by a simple rung shift within the quarter-step
lattice: the predicted $7/2$ is the unique deep-window
quarter-step choice.  This makes the tension a genuine
quantitative test rather than a knob to be adjusted.
\end{remark}

\subsection{Absolute mass predictions}

Absolute mass predictions require the SI calibration seam
(Paper~I, Sec.~7).  They are therefore less clean than the
seam-free predictions above, but they are included for
completeness.

\begin{table}[h]
\centering
\caption{Absolute neutrino mass predictions (conditional on
the SI anchor $\tau_0$ calibrated to $m_e=0.511$~MeV) vs.\
NuFIT~5.3 normal-ordering constraints.  The $-1/4$ phase
offset is a frontier input.}
\label{tab:masses}
\begin{tabular}{lccc}
\toprule
& $\nu_1$ & $\nu_2$ & $\nu_3$ \\
\midrule
Rung & $-239/4$ & $-231/4$ & $-217/4$ \\
$m_i^{\mathrm{pred}}$ (eV) & $\approx 0.0035$ & $\approx 0.0093$ &
  $\approx 0.050$ \\
\bottomrule
\end{tabular}
\end{table}

\begin{table}[h]
\centering
\caption{Neutrino observable predictions vs.\ NuFIT~5.3
(normal ordering).}
\label{tab:observables}
\begin{tabular}{lcc}
\toprule
Observable & RS prediction & NuFIT~5.3 (NO) \\
\midrule
$\Delta m^2_{21}$ &
  $\approx 7.3\times 10^{-5}\,\mathrm{eV}^2$ &
  $(7.53\pm 0.18)\times 10^{-5}$ \\
$|\Delta m^2_{31}|$ &
  $\approx 2.49\times 10^{-3}\,\mathrm{eV}^2$ &
  $(2.453\pm 0.033)\times 10^{-3}$ \\
$\Sigma m_\nu$ &
  $\approx 0.063\,\mathrm{eV}$ &
  $< 0.12\,\mathrm{eV}$ (cosmological) \\
\bottomrule
\end{tabular}
\end{table}

The predicted sum $\Sigma m_\nu\approx 0.063$~eV is well below
the current cosmological upper bound and will be testable by
next-generation surveys (CMB-S4, DESI, Euclid).  If these
surveys establish $\Sigma m_\nu > 0.08$~eV at high confidence,
the framework's absolute mass scale would be under pressure.

%=============================================================================
\section{Methods: RS provenance}
\label{sec:methods}
%=============================================================================

This section collects framework-specific details for readers
interested in the RS provenance.

\subsection{Origin of the neutrino baseline}

The baseline rung $r_{\nu_3}^{\mathrm{int}}=-54$ is the
negative sum of the five non-trivial hypercube integers:
\[
  V+E+F+\Epass+W = 8+12+6+11+17 = 54.
\]
The neutrino occupies the deepest ladder position, exhausting
the full integer budget of the cube geometry with a minus sign.
The active edge count $A=1$ is excluded because it enters the
charged-sector rung formula ($r_e=A+1=2$) separately.

This derivation is formally verified in Lean~4:
\texttt{Masses.BaselineDerivation.neutrino\_baseline\_eq}.

\subsection{Electroweak sector yardstick}

Neutrinos participate in electroweak interactions only, so
they use the electroweak sector yardstick
$A_{\mathrm{EW}} = 2^{+1}\cdot\Ecoh\cdot\phig^{55}
= 2\phig^{50}$ (using $\Ecoh=\phig^{-5}$).
The sector parameters $(\Bpow,r_0)=(+1,55)$ are from the
backbone yardstick table (Paper~I, Sec.~5).

\subsection{Lean verification}

The neutrino sector is formalized in
\texttt{IndisputableMonolith.Physics.NeutrinoSector} with the
following verified objects:
\begin{itemize}[nosep]
  \item \texttt{res\_nu3\_simp}, \texttt{res\_nu2\_simp},
    \texttt{res\_nu1\_simp}: the three fractional residues.
  \item \texttt{nu3\_frac\_pred\_bounds},
    \texttt{nu2\_frac\_pred\_bounds},
    \texttt{nu1\_frac\_pred\_bounds}: mass prediction bounds.
  \item \texttt{dm2\_21\_frac\_pred\_in\_nufit\_1sigma}:
    $\Delta m^2_{21}$ within NuFIT $1\sigma$.
  \item \texttt{dm2\_31\_frac\_pred\_in\_nufit\_2sigma}:
    $\Delta m^2_{31}$ within NuFIT $2\sigma$.
\end{itemize}

%=============================================================================
\section{Frontier items}
\label{sec:frontier}
%=============================================================================

This paper keeps the following items explicitly at the frontier.

\paragraph{Quarter-phase offset.}
The $-1/4$ phase offset that selects the rung triple from the
deep-ladder window is the single non-derived structural input
in the neutrino sector.  Its boundary status is documented as
item B-33 in the framework's derivation roadmap.  A plausible
geometric origin (the $C_4$ face symmetry of $Q_3$) exists but
has not been formalized into a closed first-principles argument.

All three seam-free predictions (ordering, $\phig^7$ ratio,
splitting ratio) are conditional on this offset.  If the offset
is wrong, the predictions are wrong.  This conditionality is
the honest price of publishing a prediction paper before the
derivation is complete.

\paragraph{Dirac vs.\ Majorana character.}
The framework treats neutrinos as Dirac particles ($Z_\nu=0$,
same $\phig$-ladder as charged fermions).  The question of
whether neutrinos are Dirac or Majorana is not resolved by the
present construction.  Neutrinoless double beta decay
experiments will address this question independently.

\paragraph{PMNS mixing matrix.}
A geometric derivation of the PMNS mixing angles from $Q_3$
path-weight combinatorics is sketched in the framework
(\texttt{Physics.PMNS.born\_mixing}) but is not mature enough
for a standalone prediction paper.

%=============================================================================
\section{Falsifiers}
\label{sec:falsifiers}
%=============================================================================

Each prediction carries a concrete failure mode.

\begin{enumerate}
  \item \textbf{Mass ordering.}
    If future experiments (JUNO, DUNE, Hyper-Kamiokande)
    establish inverted ordering at $>3\sigma$, the rung
    triple~\eqref{eq:triple} and the entire neutral-sector
    construction are refuted.

  \item \textbf{Splitting ratio.}
    If $\Rdel^{\mathrm{exp}}$ is measured to lie outside the
    interval $(30.0,\,37.6)$ at $>3\sigma$ (corresponding to
    $33.82\pm 3\times 1.3$, where $1.3$ is the framework's
    structural resolution from adjacent quarter-step
    alternatives), the edge-partition rung assignment is
    falsified.

    More precisely: if future oscillation data push the central
    value of $\Rdel$ below $\approx 32.2$ at $>3\sigma$ from
    $33.82$, the framework is refuted.  The current $1.5\sigma$
    tension is the strongest existing pressure.

  \item \textbf{$\phig^7$ squared-mass ratio.}
    If absolute neutrino masses are measured (e.g., from
    KATRIN endpoint analysis combined with oscillation data)
    and $m_3^2/m_2^2$ deviates from $\phig^7\approx 29.03$
    by more than $3\sigma$, the atmospheric--solar rung gap
    $\Delta_{32}=7/2$ is refuted.

  \item \textbf{Sum of masses.}
    If $\Sigma m_\nu$ is measured to exceed $0.08$~eV at high
    significance, the absolute mass scale predicted by the
    electroweak sector yardstick is under pressure.  If
    $\Sigma m_\nu > 0.12$~eV, the prediction
    $\Sigma m_\nu\approx 0.063$~eV is strongly disfavored.

  \item \textbf{Quarter-phase ablation.}
    If any of the alternative deep-window quarter-step choices
    ($4r_3=-219$ or $4r_3=-218$) produces a splitting ratio
    closer to experiment, the $-1/4$ phase offset is called
    into question.  Currently, neither alternative produces
    viable normal-ordering predictions.
\end{enumerate}

%=============================================================================
\section{Discussion}
\label{sec:discussion}
%=============================================================================

The neutrino sector is the framework's sharpest testing ground
for two reasons.  First, the neutral-sector simplification
($\mathrm{gap}(0)=0$) eliminates the charge-band correction
entirely, leaving a pure $\phig$-ladder.  Second, the
splitting ratio $\Rdel$ is dimensionless and depends on nothing
but $\phig$ and two cube integers ($\Epass=11$ and
$2^{D-1}=4$).  No other sector prediction achieves this level
of cleanness.

The current $1.5\sigma$ tension between the predicted
$\Rdel\approx 33.82$ and the NuFIT central value $31.6$ is
the most important open datum for the framework.  It is neither
comfortable nor fatal.  Next-generation oscillation experiments
(JUNO, expected to reach sub-percent precision on $\Delta m^2_{21}$;
DUNE, expected to reach percent-level precision on
$\Delta m^2_{31}$) will either sharpen the tension into a
refutation or relax it into consistency.  Either outcome is
scientifically productive.

The single frontier item---the $-1/4$ phase offset---is an
honest weakness.  But it is a weakness that is kept visible and
explicitly conditional, not hidden behind suggestive prose.
The framework's neutrino predictions are strong enough to be
worth publishing even in their conditional form, precisely
because they are falsifiable.

%=============================================================================
\section{Conclusion}
%=============================================================================

Given the neutral-sector quarter-step structure of the
Recognition Science mass backbone, three predictions follow
that are independent of the SI calibration seam:
\begin{align}
  &\text{Ordering:} &&m_1 < m_2 < m_3
    \quad\text{(normal)}, \label{eq:conc_ord} \\
  &\text{Squared-mass ratio:} &&m_3^2/m_2^2 = \phig^7
    \approx 29.03, \label{eq:conc_phi7} \\
  &\text{Splitting ratio:} &&\Rdel
    = \frac{\phig^{11}-1}{\phig^4-1}
    \approx 33.82. \label{eq:conc_split}
\end{align}

The splitting ratio is the sharpest near-term test.  It uses
only $\phig$ and the cube integers $\Epass=11$ and $2^{D-1}=4$.
If future data push $\Rdel$ beyond $3\sigma$ from $33.82$,
the framework is refuted.

The predictions are conditional on one frontier input: the
$-1/4$ quarter-phase offset.  This conditionality is stated
plainly because it is better science to publish a sharp
conditional prediction than to wait indefinitely for a complete
derivation that may never arrive, or to pretend the condition
has already been discharged.

\section*{Acknowledgments}

The author thanks Elshad Allahyarov for detailed manuscript
review and collaboration on the neutrino sector closure.

\begin{thebibliography}{9}

\bibitem{backbone}
J.~Washburn,
``The Recognition Science Fermion Mass Backbone at the Anchor Scale:
Structural Mass Coordinates, Sector Yardsticks, and the Explicit
Calibration Seam,''
Recognition Physics Institute preprint, March 2026.

\bibitem{leptons}
J.~Washburn,
``Charged Lepton Mass Ratios from the Combinatorics of the 3-Cube,''
Recognition Physics Institute preprint, March 2026.

\bibitem{CKM}
J.~Washburn,
``From Geometric Rungs to Standard Model Vertices: CKM Geometry and
Yukawa Matching in Recognition Science,''
Recognition Physics Institute preprint, March 2026.

\bibitem{fullderivation}
J.~Washburn,
``Complete Fermion Mass Derivation in Recognition Science,''
Recognition Physics Institute internal draft, March 2026.

\bibitem{NuFIT}
I.~Esteban, M.~C.~Gonzalez-Garcia, M.~Maltoni, T.~Schwetz, and
A.~Zhou,
``The fate of hints: updated global analysis of three-flavor
neutrino oscillations,''
JHEP \textbf{09} (2020) 178;
NuFIT~5.3 (2024), \url{http://www.nu-fit.org}.

\end{thebibliography}

\end{document}
